\documentclass[twoside]{article}

\usepackage[preprint]{aistats2027}
\usepackage[round,authoryear]{natbib}
\usepackage{amsmath,amssymb,amsthm}
\usepackage{mathtools}
\usepackage{microtype}
\usepackage{url}

\renewcommand{\bibname}{References}
\renewcommand{\bibsection}{\subsubsection*{\bibname}}
\bibliographystyle{apalike}

\newtheorem{theorem}{Theorem}
\newtheorem{lemma}[theorem]{Lemma}
\newtheorem{proposition}[theorem]{Proposition}
\newtheorem{corollary}[theorem]{Corollary}
\newcommand{\E}{\mathbb E}
\newcommand{\cH}{\mathcal H}
\newcommand{\FM}{\textnormal{\textsc{ForceMoM}}}

\begin{document}
\runningtitle{Adaptation Frontier for Heavy-Tailed Bandits}

\twocolumn[
\aistatstitle{Supplement: A Sharp Adaptation Frontier for Heavy-Tailed Bandits\\with Unknown Moments}

\aistatsauthor{}

\aistatsaddress{}
]

This supplement gives complete proofs of all results in the main paper. It
also reproves the unknown-scale lower tradeoff with a corrected constant.

\section{Algorithm and Exploration Counts}
\label{app:schedule}

Fix $K\geq2$ and $\rho\in(0,1)$. Let
\[
 m_0=\lceil16\log(2K+e)\rceil,
 \qquad t_0=Km_0.
\]
The algorithm first pulls each arm $m_0$ times in cyclic order. These pulls
are marked as exploration. For $t>t_0$, define
\[
 D_t=\max\{t_0,\lceil K^{1-\rho}t^\rho\rceil\}.
\]
If the exploration count before round $t$ is smaller than $D_t$, the
algorithm explores the next arm in cyclic order. Otherwise it exploits.

Let $F_t$ be the total number of exploration pulls through round $t$, and
let $N_i^E(t)$ be the number assigned to arm $i$. The initialization has
$F_{t_0}=D_{t_0}=t_0$.

\begin{lemma}[Exact tracking]
\label{lem:tracking}
For every $t\geq t_0$, $F_t=D_t$. Moreover,
\begin{equation}
 \left\lfloor\frac{D_t}{K}\right\rfloor
 \leq N_i^E(t)
 \leq\left\lceil\frac{D_t}{K}\right\rceil
 \label{eq:cyclic-count}
\end{equation}
for every arm $i$.
\end{lemma}

\begin{proof}
Let $g(t)=K^{1-\rho}t^\rho$. For $t\geq K$,
\[
 g'(t)=\rho K^{1-\rho}t^{\rho-1}\leq\rho<1.
\]
Hence $g(t)-g(t-1)<1$ and
$\lceil g(t)\rceil-\lceil g(t-1)\rceil\in\{0,1\}$. Taking the maximum with
the constant $t_0$ preserves this property. Since
$g(t_0)=K m_0^\rho\leq Km_0=t_0$, we have $D_{t_0}=t_0$.

Inductively, suppose $F_{t-1}=D_{t-1}$. If $D_t=D_{t-1}$, the condition for
forced exploration is false and $F_t=D_t$. If $D_t=D_{t-1}+1$, the
condition is true, so exactly one exploration pull occurs and again
$F_t=D_t$. This proves exact tracking. Cyclic assignment makes the arm
counts differ by at most one, which gives \eqref{eq:cyclic-count}.
\end{proof}

Before any exploitation decision at $t>t_0$, every arm therefore has at
least
\begin{equation}
 n_t:=\left\lfloor D_{t-1}/K\right\rfloor
 \geq\left\lfloor K^{-\rho}(t-1)^\rho\right\rfloor.
 \label{eq:n-lower-floor}
\end{equation}
Since $t_0/K=m_0\geq31$, the quantity inside the last floor is at least
$m_0^\rho\geq1$. Since $\lfloor x\rfloor\geq x/2$ for every $x\geq1$,
we have uniformly for $t>t_0$ that
\begin{equation}
 n_t\geq\frac12K^{-\rho}(t-1)^\rho.
 \label{eq:n-lower}
\end{equation}

The exploration regret admits both uniform and instance-specific bounds.
Jensen's inequality implies $\Delta_i\leq2u^{1/p}$, so
\begin{equation}
 R_T^E\leq2u^{1/p}D_T
 \leq2u^{1/p}\{Km_0+K^{1-\rho}T^\rho+1\}.
 \label{eq:explore-uniform}
\end{equation}
On a fixed instance, cyclic assignment and Lemma~\ref{lem:tracking} give
\begin{equation}
 \lim_{T\to\infty}\frac{R_T^E}{T^\rho}
 =K^{-\rho}\sum_{i=1}^K\Delta_i.
 \label{eq:explore-instance}
\end{equation}

\section{Median-of-Means Bounds}
\label{app:mom}

We prove the estimator lemma uniformly for all $p\in(1,2]$. Let
$X_1,\ldots,X_n$ be i.i.d., let $\mu=\E X$, and assume
$\E|X|^p\leq u$. Let $q\leq n$ be odd and put $k=(q+1)/2$ and
$m=\lfloor n/q\rfloor$. Form $q$ block averages $\bar X_1,\ldots,\bar X_q$
from $qm$ samples and let $M_{n,q}$ be their median.

\begin{lemma}[One block]
\label{lem:block}
For every block $j$,
\begin{equation}
 \E|\bar X_j-\mu|^p
 \leq16u(q/n)^{p-1}.
 \label{eq:block-moment}
\end{equation}
\end{lemma}

\begin{proof}
By Jensen, $|\mu|^p\leq u$. Convexity gives
\[
 \E|X-\mu|^p
 \leq2^{p-1}\{\E|X|^p+|\mu|^p\}
 \leq2^p u.
\]
The von Bahr--Esseen inequality
\citep{vonbahr1965inequalities} yields
\begin{align*}
 \E|\bar X_j-\mu|^p
 &\leq m^{-p}\,2m\E|X-\mu|^p\\
 &\leq2^{p+1}u m^{1-p}.
\end{align*}
Since $n/q\geq1$, $m=\lfloor n/q\rfloor\geq n/(2q)$. Also
$2^{p+1}2^{p-1}=2^{2p}\leq16$. Substitution proves
\eqref{eq:block-moment}.
\end{proof}

\begin{lemma}[MoM tail and mean]
\label{lem:mom-full}
For every $x>0$,
\begin{equation}
 \Pr\{|M_{n,q}-\mu|>x\}
 \leq\min\left\{1,
 \left(\frac{64u(q/n)^{p-1}}{x^p}\right)^k\right\}.
 \label{eq:mom-tail-app}
\end{equation}
If $q\geq3$, then
\begin{equation}
 \E|M_{n,q}-\mu|
 \leq128u^{1/p}(q/n)^{(p-1)/p}.
 \label{eq:mom-mean-app}
\end{equation}
\end{lemma}

\begin{proof}
By Lemma~\ref{lem:block} and Markov's inequality, the probability that one
block differs from $\mu$ by more than $x$ is at most
\[
 \theta_x=\min\{1,16u(q/n)^{p-1}x^{-p}\}.
\]
If the median differs by more than $x$, at least $k$ blocks do. Their
independence and a union bound over subsets of size $k$ give
\[
 \Pr\{|M_{n,q}-\mu|>x\}
 \leq2^q\theta_x^k.
\]
Since $q=2k-1$, $2^q\leq4^k$. This proves
\eqref{eq:mom-tail-app}.

Set
\[
 x_0=64^{1/p}u^{1/p}(q/n)^{(p-1)/p}.
\]
Integrating \eqref{eq:mom-tail-app} gives
\begin{align*}
 \E|M_{n,q}-\mu|
 &\leq x_0+\int_{x_0}^\infty(x_0/x)^{pk}\,dx\\
 &=x_0\left(1+\frac1{pk-1}\right).
\end{align*}
For $q\geq3$, $k\geq2$ and $pk-1>1$. Hence the last display is at most
$2x_0\leq128u^{1/p}(q/n)^{(p-1)/p}$.
\end{proof}

\begin{lemma}[Maximum over arms]
\label{lem:max-full}
Suppose $K$ arm estimators satisfy \eqref{eq:mom-tail-app} with arm-specific
$q_i,n_i,k_i$, where $q_i\leq q$, $n_i\geq n$, and
$k_i\geq\log(2K)$. Then
\begin{equation}
 \E\max_{i\leq K}|M_i-\mu_i|
 \leq128e\,u^{1/p}(q/n)^{(p-1)/p}.
 \label{eq:max-app}
\end{equation}
No independence across arms is required.
\end{lemma}

\begin{proof}
Let $k_0=\min_i k_i$ and
$A=64u(q/n)^{p-1}$. The union bound and
\eqref{eq:mom-tail-app} imply
\[
 \Pr\{\max_i|M_i-\mu_i|>x\}
 \leq\min\{1,K(A/x^p)^{k_0}\}
\]
whenever $x^p\geq A$. Let
$x_0=A^{1/p}K^{1/(pk_0)}$. Tail integration as in the previous proof gives
an expectation at most $2x_0$. Since $k_0\geq\log(2K)$,
\[
 K^{1/(pk_0)}
 \leq\exp\{\log K/\log(2K)\}\leq e.
\]
Also $2A^{1/p}\leq128u^{1/p}(q/n)^{(p-1)/p}$.
\end{proof}

For the algorithm, let $\ell:\mathbb N\to[1,\infty)$ be nondecreasing with
$\ell(n)\to\infty$, $\ell(n+1)-\ell(n)\leq1$, and
$\log\ell(n)=o(\log n)$. Define
\begin{align*}
 q(n)=\max\{r:\;&r\leq n,\ r\leq
 2\lceil\log(2K)+\ell(n)\rceil+1,\\
 &r\text{ is odd}\}.
\end{align*}
The assumed schedule is nondecreasing and satisfies
$\ell(n+1)-\ell(n)\leq1$. Hence $q(n+1)-q(n)\leq2$. This is the only
regularity property of the block schedule needed below.
For $n\geq m_0$, both arguments of the minimum exceed
$2\log(2K)+1$. Thus $q(n)\geq2\log(2K)-1$, and
$k(n)=(q(n)+1)/2\geq\log(2K)$. Also
\begin{equation}
 q(n)\leq3+2\log(2K)+2\ell(n).
 \label{eq:q-upper-app}
\end{equation}

\section{Proof of the Finite-Time Upper Bound}
\label{app:upper}

Let $J_t$ be the arm chosen on an exploitation round and fix an optimal arm
$i_*$. Since $M_{J_t}\geq M_{i_*}$,
\begin{align}
 \Delta_{J_t}
 &=\mu_*-\mu_{J_t}\nonumber\\
 &\leq|M_{i_*}-\mu_*|+|M_{J_t}-\mu_{J_t}|\nonumber\\
 &\leq2\max_i|M_i-\mu_i|.
 \label{eq:leader-app}
\end{align}
The exploration counts of any two arms differ by at most one. Applying
Lemma~\ref{lem:max-full}, \eqref{eq:n-lower}, and
\eqref{eq:q-upper-app}, we obtain for all exploitation rounds outside a
finite initial interval
\begin{align}
 \E\Delta_{J_t}
 &\leq C u^{1/p}
 \left(\frac{1+\log(2K)+\ell(t)}{K^{-\rho}t^\rho}
 \right)^{(p-1)/p}\nonumber\\
 &=C u^{1/p}K^{\rho(p-1)/p}
 L_t^{(p-1)/p}t^{-\rho(p-1)/p},
 \label{eq:instant-app}
\end{align}
where $L_t=1+\log(2K)+\ell(t)$.

Put $a=\rho(p-1)/p$. Since $p\leq2$ and $\rho<1$, $a<1/2$. Monotonicity of
$L_t$ and integral comparison give
\begin{equation}
 \sum_{t=1}^T L_t^{(p-1)/p}t^{-a}
 \leq2L_T^{(p-1)/p}T^{1-a}.
 \label{eq:sum-app}
\end{equation}
Combining \eqref{eq:explore-uniform}, \eqref{eq:instant-app}, and
\eqref{eq:sum-app}, and enlarging the universal constant to cover
$T\leq t_0$, proves
\begin{align*}
 R_T\leq C u^{1/p}\big\{&K\log(eK)+K^{1-\rho}T^\rho\\
 &+K^{\rho(p-1)/p}T^{1-\rho(p-1)/p}
 L_T^{(p-1)/p}\big\}.
\end{align*}
This is Theorem 2 of the main paper.

If $\rho\leq2/3$, then for every $p\in(1,2]$,
\[
 1-\rho(p-1)/p-\rho
 =1-\rho(2p-1)/p
 \geq1-3\rho/2\geq0.
\]
Thus the exploitation term has the larger or equal exponent.

\section{Proof of the Fixed-Instance Bound}
\label{app:instance}

Fix an instance in $\cH_{p,u}$ and a suboptimal arm $i$ with
$\Delta_i>0$. If $i$ is selected on an exploitation round, then
\[
 M_i\geq M_{i_*}.
\]
This implies either $M_i-\mu_i\geq\Delta_i/2$ or
$\mu_*-M_{i_*}\geq\Delta_i/2$. Let $\underline q_t$ and
$\overline q_t$ be the smaller and larger block counts among the two
estimates. The arm exploration counts differ by at most one, so
$\overline q_t\leq\underline q_t+2$. By Lemma~\ref{lem:mom-full}, for a
constant $C$ and all sufficiently large $t$,
\begin{equation}
 \Pr\{J_t=i\}
 \leq2\left[
 \frac{Cu \overline q_t^{p-1}}
 {\Delta_i^p n_t^{p-1}}
 \right]^{(\underline q_t+1)/2}.
 \label{eq:mistake-app}
\end{equation}
Here $n_t$ is the smaller sample count.

We verify summability. Since $\log\ell(n)=o(\log n)$,
\[
 \log \overline q_t=o(\log n_t).
\]
Hence, for every fixed $p,u,\Delta_i$, the logarithm of the quantity inside
the brackets in \eqref{eq:mistake-app} is eventually at most
\[
 -\frac{p-1}{2}\log n_t.
\]
By \eqref{eq:n-lower}, it is then at most
$-(p-1)\rho\log t/4$ after increasing the instance-dependent threshold.
Since $\underline q_t\to\infty$, eventually
\[
 \frac{(p-1)\rho(\underline q_t+1)}8\geq3.
\]
Therefore $\Pr\{J_t=i\}\leq2t^{-3}$ eventually. Summing over time and over
the finitely many suboptimal arms proves
\[
 \sum_{t=1}^\infty
 \Pr\{t\text{ is exploit and }J_t\text{ is suboptimal}\}<\infty.
\]
The expected total regret due to exploitation is finite. Equation
\eqref{eq:explore-instance} now proves
\[
 \limsup_{T\to\infty}\frac{R_T}{T^\rho}
 \leq K^{-\rho}\sum_i\Delta_i.
\]

\section{The Unknown-Scale Lower Tradeoff}
\label{app:lower}

For completeness, we prove the exponent lower bound used in the main paper.
The argument follows \citet{genalti2025lower}, which adapts
\citet{hadiji2023range}. Our statement uses a $p$-dependent constant. This is
all that is needed for the polynomial frontier.

\begin{theorem}[Lower tradeoff]
\label{thm:lower-app}
Fix $p\in(1,2]$. Suppose a policy does not know $u$ and, for a function
$\Phi(T)=o(T)$, satisfies
\begin{equation}
 R_T(\eta)\leq v^{1/p}\Phi(T)
 \label{eq:moment-free-app}
\end{equation}
for every $v>0$, every $\eta\in\cH_{p,v}$, and every $T$. Then every fixed
instance $\nu$ with at least one suboptimal arm satisfies
\begin{equation}
 \liminf_{T\to\infty}
 \frac{R_T(\nu)}{(T/\Phi(T))^{p/(p-1)}}
 \geq c_p\sum_{i:\Delta_i>0}\Delta_i,
 \label{eq:lower-app}
\end{equation}
where one may take
\begin{equation}
 c_p=\frac{1}{2\,8^{p/(p-1)}}.
 \label{eq:cp}
\end{equation}
\end{theorem}

\begin{proof}
Fix a suboptimal arm $a$ of $\nu$, with law $\nu_a$, mean $\mu_a$, and gap
$\Delta_a>0$. Let $v$ be any finite common upper bound on the arms' $p$th
absolute moments. For $\beta\in(0,1/2)$, define an alternative law
\begin{equation}
 \nu'_{a,\beta}
 =(1-\beta)\nu_a+
 \beta\delta_{\mu_a+2\Delta_a/\beta}.
 \label{eq:alternative}
\end{equation}
Its mean is $\mu_a+2\Delta_a=\mu_*+\Delta_a$, so arm $a$ is uniquely better
than every original optimal arm by at least $\Delta_a$. Let $v'_\beta$ be a
common $p$th-moment bound for the alternative bandit. We may take
\begin{equation}
 v'_\beta
 =v+(1-\beta)\E_{\nu_a}|X|^p
 +\beta|\mu_a+2\Delta_a/\beta|^p.
 \label{eq:vprime}
\end{equation}
The redundant $v$ term safely covers all unchanged arms.

Write $N_a(T)$ for the number of pulls of arm $a$, and let $P$ and $P'$ be
the history laws under the original and alternative instances. Mixture
domination gives
\[
 \mathrm{KL}(\nu_a,\nu'_{a,\beta})
 \leq\log\frac1{1-\beta}\leq2\beta\log2.
\]
The standard bandit change-of-measure identity and data processing give
\begin{equation}
 \operatorname{kl}\left(
 \frac{\E_PN_a(T)}T,
 \frac{\E_{P'}N_a(T)}T\right)
 \leq2\beta\log2\,\E_PN_a(T).
 \label{eq:data-processing}
\end{equation}
For $x,y\in[0,1]$, binary relative entropy obeys
\begin{equation}
 \operatorname{kl}(x,y)
 \geq(1-x)\log\frac1{1-y}-\log2.
 \label{eq:binary-kl}
\end{equation}

The moment-free guarantee implies
\begin{align}
 \E_PN_a(T)&\leq \frac{v^{1/p}\Phi(T)}{\Delta_a},
 \label{eq:count-original}\\
 T-\E_{P'}N_a(T)&\leq
 \frac{(v'_\beta)^{1/p}\Phi(T)}{\Delta_a}.
 \label{eq:count-alt}
\end{align}
Substituting \eqref{eq:count-original} and \eqref{eq:count-alt} into
\eqref{eq:data-processing} through \eqref{eq:binary-kl} yields
\begin{align}
 &\left(1-\frac{v^{1/p}\Phi(T)}{T\Delta_a}\right)
 \log\left(\frac{T\Delta_a}{(v'_\beta)^{1/p}\Phi(T)}\right)-\log2
 \nonumber\\
 &\hspace{35mm}\leq2\beta\log2\,\E_PN_a(T).
 \label{eq:key-lower}
\end{align}

Put $r=p/(p-1)$ and choose
\begin{equation}
 \alpha=8^{-r},
 \qquad
 \beta_T=\alpha^{-1}\left(\frac{\Phi(T)}T\right)^r.
 \label{eq:beta-choice}
\end{equation}
Since $\Phi(T)=o(T)$, $\beta_T\to0$. Equation \eqref{eq:vprime} gives
\begin{equation}
 (v'_{\beta_T})^{1/p}\Phi(T)
 =2\Delta_a\alpha^{(p-1)/p}T\{1+o(1)\}.
 \label{eq:vprime-limit}
\end{equation}
The right side equals $\Delta_a T\{1+o(1)\}/4$.
The first factor on the left of \eqref{eq:key-lower} tends to one. The
logarithm tends to $\log4$, so the entire left side tends to $\log2$.
Using \eqref{eq:beta-choice} on the right gives
\begin{align*}
 \liminf_{T\to\infty}
 \frac{\E_PN_a(T)}{(T/\Phi(T))^r}
 &\geq\frac{\alpha}{2\log2}\log2\\
 &=\frac1{2\,8^r}=c_p.
\end{align*}
Multiplying by $\Delta_a$ and summing this inequality over all suboptimal
arms proves \eqref{eq:lower-app}.
\end{proof}

If $\Phi(T)=T^{a+o(1)}$ and a fixed-instance rate is
$T^{b+o(1)}$, Theorem~\ref{thm:lower-app} implies
\[
 b\geq\frac{p}{p-1}(1-a),
\]
or equivalently
\[
 b+\frac{p}{p-1}a\geq\frac{p}{p-1}.
\]
For the algorithmic exponent
$a=1-\rho(p-1)/p$, the right side is exactly $\rho$. This completes the
proof of the simultaneous sharp frontier.

\bibliography{references}

\end{document}
